
\section{Moltiplicazione di polinomi - Karatsuba}
\subsubsection*{Definizione formale di polinomio e degree bound}
Un polinomio \(P(x)\) di grado \(N-1\) (con \(N\) termini) è definito come:
\[P(x) \; := \; \sum_{i=0}^{N-1} \, a_i \, x^i \qquad\qquad \text{con} \;\; a_0, a_1, \dotsc \; a_{N-1} \; \text{fissati}\]
Si osserva che un polinomio può essere visto anche come un numero espresso in base \(x\) con \(0\leq a_i < x\) secondo la
notazione posizionale.

Il degree bound \(\alpha\) è un limite superiore al grado di un polinomio e indica l'insieme di tutti i polinomi di
grado inferiore o uguale ad \(\alpha\), ovvero quelli con \(N - 1 \leq \alpha\).

\subsection{Prodotto tra polinomi usando la definizione}
\subsubsection*{Tecnica risolutiva}
Dati due polinomi \(P(x)\) e \(Q(x)\) di grado \(N - 1\), il loro prodotto \(P(x) \cdot Q(x)\) è anch'esso un poliomio
di grado \(2N - 2\), come illustrato di seguito.
\[P(x) \; := \; \sum_{i=0}^{N-1} \, a_i \, x^i \qquad\quad Q(x) \; := \; \sum_{i=0}^{N-1} \, b_i \, x^i \qquad\quad P(x) \cdot Q(x) \; = \; \sum_{i=0}^{N-1} \sum_{j=0}^{N-1} \, a_i \, b_j \, x^{i+j} \; = \; \sum_{k=0}^{2N-2} \, c_k \, x^k\]

\subsubsection*{Analisi della complessità}
Per calcolare il prodotto usando la definizione di prodotto tra polinomi sono necessarie:
\begin{itemize}
	\item \(N^2\) moltiplicazioni tra i coefficienti \(a_i\) e \(b_j\)
	\item \((N-1)^2 = N^2 - 1 - (2N - 2)\) addizioni tra i prodotti \(a_i \cdot b_j\) per ottenere i coefficienti \(c_k\)
\end{itemize}
Complessivamente il prodotto usando la definizione ha complessità \(\Theta(N^2)\).

\subsection{Prodotto tra polinomi con approccio divide and conquer}
\subsubsection*{Tecnica risolutiva}
I polinomi \(P(x)\) e \(Q(x)\), di grado \(N-1\), possono essere scomposti in polinomi più semplici di grado \(< N/2\)
e il prodotto tra \(P(x)\) e \(Q(x)\) è esprimibile in funzione del prodotto di questi polinomi più semplici, ovvero:
\[P(x) = A(x) + B(x) \cdot x^{N/2} \qquad\qquad Q(x) = C(x) + D(x) \cdot x^{N/2}\]
\[P(x) \cdot Q(x) = B(x) \, D(x) \cdot x^N + (A(x) \, D(x) + B(x) \, C(x)) \cdot x^{N/2} + A(x) \, C(x)\]

\subsubsection*{Equazione alle ricorrenze}
L'equazione alle ricorrenze per questo approccio è la seguente:
\[T(N) \; = \; \begin{cases}
	\; 1 & \text{se } N = 1 \\[3pt]
	\; 4 \cdot T(N/2) + \beta_{std} \; N & \text{se } N > 1
\end{cases}\]
\begin{itemize}
	\item il caso base è rappresentato dal prodotto di due polinomi di grado 0 che richiede un solo prodotto
	\item il passo ricorsivo prevede 4 chiamate ricorsive per calcolare i quattro prodotti \(AC\), \(AD\), \(BC\) e \(BD\)
	tra polinomi di grado \(< N/2\)
	\item il costo per ogni nodo interno è dato dalle somme dei coefficienti dei termini di grado uguale, che è proporzionale
	alla taglia \(N\)
\end{itemize}

\subsubsection*{Analisi della complessità}
Dal formulone si ottiene che la complessità rimane sempre \(\Theta(N^2)\), come evidenziato di seguito:
\begin{align*}
	T(N) \;\; &= \sum_{l=0}^{(\log_2 N) - 1} \!\! \left( 4^l \cdot \beta_{std} \; \frac{N}{2^l} \right) + 4^{\log_2 N} \cdot T_0(n) \; = \; \beta_{std} \; N \cdot \! \sum_{l=0}^{(\log_2 N) - 1} \!\! \left( 2^{2l}/2^l \right) + N^{\log_2 4} \; = \\[3pt]
	&= \;\; \beta_{std} \; N \cdot \sum_{l=0}^{(\log_2 N) - 1} \!\! 2^l + N^2 \; = \; \beta_{std} \; N \cdot \frac{2^{\log_2 N} - 1}{2 - 1} + N^2 \; = \\[3pt]
	&= \, \underbrace{\, \beta_{std} \; (N^2 - N) \,}_{\begin{array}{c} \text{\scriptsize costo delle somme} \\[-3pt] \text{ \scriptsize nei nodi interni} \end{array}}
	\;\; + \!\!\!\!\!\! \underbrace{\;\;\;N^2\;\;}_{\begin{array}{c} \text{\scriptsize costo dei prodotti} \\[-3pt] \text{\scriptsize nelle foglie} \end{array}} \!\!\!\!\!\! = \;\; \Theta \, (N^2)
\end{align*}

\subsection{Prodotto di polinomi con l'algoritmo di Karatsuba (1962)}
\subsubsection*{Tecnica risolutiva}
Karatsuba osserva che è possibile ricavare la somma dei due prodotti \(AD + BC\) riutilizzando i due prodotti \(AC\) e \(BD\)
già precedentemente calcolati come segue:
\[(A + B) \cdot (C + D) = AC + (AD + BC) + BD \implies AD + BC = (A + B) \cdot (C + D) - AC - BD\]
In questo modo il numero di prodotti da calcolare ricorsivamente si riduce da 4 a 3, con l'effetto di aumentare lievemente
il costo della fase di divide dovuto al calcolo delle somme \((A + B)\) e \((C + D)\).

\subsubsection*{Equazione alle ricorrenze ed analisi della complessità}
L'equazione alle ricorrenze per l'algoritmo di Karatsuba è la seguente:
\[T(N) \; = \; \begin{cases}
	\; 1 & \text{se } N = 1 \\[3pt]
	\; 3 \cdot T(N/2) + \beta_{kz} \; N & \text{se } N > 1
\end{cases}\]

\noindent
Dal formulone si ottiene che la complessità è \(\Theta(N^{\log_2 3}) \approx \Theta(N^{1.58})\), come evidenziato di seguito:
\begin{align*}
	T(N) \;\; &= \sum_{l=0}^{(\log_2 N) - 1} \!\! \left( 3^l \cdot \beta_{kz} \; N / 2^l \right) + 3^{\log_2 N} \cdot T(1) \; = \; \beta_{kz} \; N \cdot \! \sum_{l=0}^{(\log_2 N) - 1} \!\! \left( (3/2)^l \right) + N^{\log_2 3} \; = \\[3pt]
	&= \;\; \beta_{kz} \; N \cdot \frac{(3/2)^{\log_2 N} - 1}{(3/2) - 1} + N^{\log_2 3} \; = \; 2 \beta_{kz} \; N \cdot (N^{\log_2 3 - \log_2 2} - 1) + N^{\log_2 3} \; = \\[3pt]
	&= \;\; \underbrace{\, 2 \beta_{kz} \; (N^{\log_2 3} - N) \,}_{\begin{array}{c} \text{\scriptsize costo delle somme} \\[-3pt] \text{\scriptsize nei nodi interni} \end{array}}
	\;\; + \!\!\!\!\!\! \underbrace{\;\;N^{\log_2 3}\;\;}_{\begin{array}{c} \text{\scriptsize costo dei prodotti} \\[-3pt] \text{\scriptsize nelle foglie} \end{array}} \!\!\!\!\!\! = \;\; \Theta \, (N^{\log_2 3}) \; \approx \; \Theta \, (N^{1.58})
\end{align*}

\subsubsection*{Considerazioni sull'algoritmo di Karatsuba}
\begin{itemize}
	\item per valori di \(N\) piccoli, l'algoritmo standard (D\&C) rimane ancora più efficiente di quello di Karatsuba,
	in quanto il costo aggiuntivo della fase di divide è maggiore del risparmio ottenuto dalla chiamata ricorsiva in meno
	\item per valori di \(N\) molto grandi, è possibile migliorare ulteriormente l'efficienza asintotica nel calcolo del
	prodotto di polinomi utilizzando ad esempio l'algoritmo di Toom-Cook
	\item la vera novità dell'algoritmo di Karatsuba è stata quella di trovare il primo algoritmo sub-quadratico per
	risolvere un problema che tutti i più grandi matematici ritenevano impossibile da risolvere in meno di \(\Theta(N^2)\)
\end{itemize}
